\documentclass[10pt,a4paper]{amsart}
\usepackage[margin=19mm]{geometry}
\usepackage{amsmath,amssymb,mathtools}
\usepackage[scr=rsfs,cal=euler]{mathalfa}
\usepackage{xfrac}
\usepackage[unicode,hidelinks,bookmarks=false]{hyperref}
\newcommand{\ie}{i.e.\ }
\newcommand{\cf}{cf.\ }
\newcommand{\resp}{respectively\ }
\newcommand{\Subsection}{\subsection}
\providecommand{\nobreakdash}{\nobreak\hbox{-}}
\setlength{\emergencystretch}{2em}
\multlinegap0pt
\usepackage{bm}
\usepackage{bbm}
\usepackage{dsfont}
\usepackage{xspace}
\usepackage{cancel}
\usepackage{mathtools}
\usepackage{amsthm}
\makeatletter
\@ifundefined{proposition}{\newtheorem{proposition}{Proposition}}{}
\makeatother
\newcommand{\R}{\mathbb{R}}  % real numbers 
\newcommand{\ep}{\varepsilon}  
\newcommand{\pp}{p}
\renewcommand{\AA}{A}  % shape of interest
\renewcommand{\d}{\mathrm{d}}  % for integral e.g. \d x
\title[PDE methods for extracting normal vector fields and distance]{PDE methods for extracting normal vector fields and distance functions of shapes}
\author{Takahiro Hasebe, Jun Masamune, Hiroshi Teramoto, Takayuki Yamada}
\date{Source: arXiv:2506.19323v1 (CC BY 4.0)}
\begin{document}
\maketitle

\noindent\emph{Source and licence.} PDE methods for extracting normal vector fields and distance functions of shapes. Version arXiv:2506.19323v1 (CC BY 4.0), distributed under the Creative Commons Attribution 4.0 International licence, as recorded in the arXiv licence metadata for this exact version and shown on the abstract page https://arxiv.org/abs/2506.19323v1. Full rights notice: licenses/arxiv-2506.19323-CC-BY-4.0.txt.\par\smallskip

\noindent\textit{Editorial scope.} Counted unit: the Proposition whose source label is \texttt{Thm:piecewise} in the article (source0.tex lines 697--710, source label \texttt{Thm:piecewise}), delivered together with the article's own complete original proof of it (source0.tex lines 714--730, one \texttt{proof} environment of 17 lines). No mathematics is written, repaired or re-derived: statement and proof are the article's own text line for line. The proof is detached from the statement in the source: the article states the result at lines 697--710 and prints its proof at lines 714--730, and the proof environment's own header names the statement, so the association is the article's own. Required context retained uncounted: eq:asymptotic\_normal (lines 326--328); eq:E3 (lines 385--388); eq:n\_corner (lines 571--573). Every definition, display and label of those lines is retained unchanged. Omitted from this package and not redistributed: 1--325, 329--384, 389--570, 574--696, 711--713, 731--1153. Nothing inside a retained excerpt is omitted or altered: every retained body is delivered exactly as the article's lines are written, so no omission receipt of any kind accompanies this package. Citations: the retained excerpts carry no citation command at all, so no bibliography entry of the article is transcribed and no reference is redistributed. Reference adaptation: none was needed and none was made; every reference inside the delivered unit resolves inside it, so no unresolved reference or citation is produced. Printed slips kept literal: no printed word, symbol, hypothesis or number of the counted statement or of its proof was altered. Layout and class: the article's own class is \texttt{svmult} with packages including makeidx, graphicx, multicol, footmisc, newtxtext, newtxmath, subfigure, tikz, color; the wrapper is the lane's pinned \texttt{amsart} editorial wrapper at 10pt on A4, which restates the theorem-like environments the retained text needs and adds the article's own macro definitions that the retained text needs (\texttt{\textbackslash AA}, \texttt{\textbackslash R}, \texttt{\textbackslash d}, \texttt{\textbackslash ep}, \texttt{\textbackslash pp}), with extra wrapper packages (bm, bbm, dsfont, xspace, cancel, mathtools). The delivered numbering is the unit's own. Assets: no retained span contains a figure environment, a graphics-inclusion command, a PDF-page inclusion, a table or a TikZ picture; no image file is copied into this package, the package declares no asset dependency and no image file is redistributed with it. Licence: version arXiv:2506.19323v1 of this article is distributed under the Creative Commons Attribution 4.0 International licence, as recorded in the arXiv licence metadata for this exact version and shown on the abstract page https://arxiv.org/abs/2506.19323v1. A full rights notice is delivered at licenses/arxiv-2506.19323-CC-BY-4.0.txt. Characters: every retained excerpt is pure ASCII, so the delivered engine, XeLaTeX with its default Latin Modern text font, reports no missing character.
\begin{equation}\label{eq:asymptotic_normal}
\nabla_x h(t,x) = -  t^{-\frac1{2}}n(x) (\xi + o(1)) \quad  \text{as}\quad t\to 0^+ 
\end{equation}

\begin{align}
\|\ep_2(t)\| &\le \left(\sup_{z \in \R^N  }\|z\| \!\cdot\!  |\kappa'(\|z\|^2)|\right) |B\Delta B'|  =O( t^\frac{\alpha}{2}R^{N+\alpha}), \label{eq:E2} \\ 
\|\ep_3(t)\| &\le \int_{\{z\in\R^N:\,\|z\|>R\}} \|z\|\!\cdot\!  |\kappa'(\|z\|^2)|\,\d z =o(1). \label{eq:E3}
\end{align}

\begin{equation}\label{eq:n_corner}
n(\pp):=\frac{n_1(\pp)+n_2(\pp)}{\|n_1(\pp)+n_2(\pp)\|}. 
\end{equation}

\begin{proposition}\label{Thm:piecewise} 
%Let $N=2, \alpha \in (0,1]$ and $\pp\in \partial\AA$ be a $C^{1,\alpha}$-corner of $\AA$. Suppose that $\AA-p$ can locally be written as \eqref{eq:A}.  
%Let $n(\pp)$ be the vector defined in \eqref{eq:n_corner}. 
Let $N=2$ and $\pp$ be a $C^{1,\alpha}$-corner of $\AA$. Let $n(\pp)$ be defined as \eqref{eq:n_corner}. We define
$$
\xi = \xi_\pp = 2\sin \frac{\phi_\pp}{2} \int_0^\infty \kappa(r^2)\,\d r
$$
 where $\phi_p\in (0,2\pi)$ is the inner angle of the curve $\partial \AA$ at $\pp$. Then formula \eqref{eq:asymptotic_normal} holds true at $x=\pp$ in the sense of pointwise convergence. 
 
% If $\pp\in\partial\AA$ is a $C^{1,\alpha}$-cusp then the following formula holds: 
% \[
%\nabla_x h(t,x) = \frac{o(1)}{\sqrt{t}}, \qquad t\to0^+.  
 %\]
 \end{proposition}

\begin{proof}[Proof of Proposition \ref{Thm:piecewise}] 
To compute the integral $\int_{C-\pp} z \kappa'(\|z  \|^2) \,\d z$, let $z= - u n(\pp) + v n_{\perp}(\pp)$ be the change of variables to $(u,v)$, where $n_\perp(\pp)$ is a unit vector orthogonal to $n(\pp)$. Let $\widehat C=\{(u,v)\in \R^2\setminus\{0\}: \arg (u+i v) \in (-\phi_\pp/2,\phi_\pp/2)\}$ which is nothing but the shifted cone $C - \pp$ in the $(u,v)$-coordinate. By the symmetry of $\widehat C$ with respect to the $u$-axis and oddness of functions, 
\begin{eqnarray*}
\int_{\substack{\|z\|\leq R\\ z\in C -\pp}} z \kappa'(\|z\|^2) \,\d z
&=& \int_{\substack{\|(u,v)\|\leq R\\ (u,v)\in \widehat C}} (-u n(\pp) + v n_{\perp}(\pp)) \kappa'(u^2+v^2) \,\d u \d v \\
&=& -n(\pp) \int_{\substack{\|(u,v)\|\leq R\\ (u,v)\in \widehat C}} u \kappa'(u^2+v^2) \,\d u \d v \\
&=& -n(\pp) \int_{\widehat C} u \kappa'(u^2+v^2) \,\d u \d v +\widetilde \ep_3(t),  
\end{eqnarray*}
where $\widetilde \ep_3(t)$ satisfies the same bound \eqref{eq:E3} used for $\ep_3(t)$. 
To finish the calculations, going to the polar coordinate amounts to 
\begin{align}
-\int_{ \widehat C} u \kappa'(u^2+v^2) \,\d u \d v \notag
&= -\int_{0}^\infty \,\d r \int_{-\phi_\pp/2}^{\phi_\pp/2}\,\d\theta \cos \theta \kappa'(r^2) r^2 \notag \\ 
&= \sin \frac{\phi_\pp}{2} \int_0^\infty \kappa(r^2)\,\d r.     \tag*{\qedhere}
\end{align}
%The case of a cusp point $\pp$ requires a single modification of the proof above. In case $\phi_\pp=0$, the volume $|B|$ is  of order $o(1)$ and so $B_{\phi_\pp}'$ can be replaced with the empty set; then the left hand side of \eqref{eq:cone} itself is of order $o(1)$.  In case $\phi_\pp=2\pi$, one can take $H_{\phi_\pp}:=\R^N$ and then the integral in the right hand side of \eqref{eq:cone} vanishes by symmetry. 
\end{proof}

\end{document}
